% Separate title page, required by Synthese for double-anonymous submission.
% Submitted alongside the anonymized manuscript (dpe.tex built with \anontrue).
\documentclass[11pt]{article}
\usepackage[letterpaper,margin=1in]{geometry}
\usepackage{lmodern}
\usepackage[T1]{fontenc}
\usepackage[utf8]{inputenc}
\usepackage[hidelinks]{hyperref}
\pagestyle{empty}

\begin{document}

\begin{center}
{\Large\bfseries Discovery Philosophy Engineering.\\[2pt]
Making Philosophical Claims Experimentally Interrogable}
\end{center}

\vspace{1.5em}

\noindent\textbf{Author.} Andrew H. Bond

\noindent\textbf{Affiliation.} Department of Computer Engineering, San Jos\'e State
University, San Jos\'e, California, United States

\noindent\textbf{ORCID.} \url{https://orcid.org/0009-0003-2599-6158}

\noindent\textbf{Corresponding author.} Andrew H. Bond,
\href{mailto:andrew.bond@sjsu.edu}{andrew.bond@sjsu.edu}

\vspace{1.5em}

\noindent\textbf{Declarations.}

\vspace{0.5em}

\noindent\emph{Funding.} No funding was received for this work.

\noindent\emph{Competing interests.} The author declares no competing interests.

\noindent\emph{Data and code availability.} Every number in Sections 7.2, 7.5, and
7.6 is read from a sealed record in a public repository. The registrations, run
logs, graders, and result files are at
\url{https://github.com/ahb-sjsu/geometric-evaluation-theory} and
\url{https://github.com/ahb-sjsu/observation-theory-campaigns}. The
specifications and the checker described in Section 7.6 are at
\url{https://github.com/ahb-sjsu/philosophy-engineering}. The source of this
manuscript is at
\url{https://github.com/ahb-sjsu/discovery-philosophy-engineering}.

\noindent\emph{Use of generative artificial intelligence.} The author used a large
language model assistant for source preparation, for the prior-art search reported
in Section 2, for drafting Sections 7.5 and 7.6 from the sealed campaign records,
and for editing. All claims, verdicts, and numbers were checked against the records
cited in each section. The author is responsible for the content.

\end{document}
